\documentclass[11pt]{article}
\usepackage{amsmath, amssymb}

\title{CSE 400: Fundamentals of Probability in Computing\\
Lecture Scribe: Inequalities and Tail Bounds}
\date{}

\begin{document}

\maketitle

\section*{1. Tail and Motivation of Tail Bounds}

\subsection*{Definition of Tail}

The tail of a random variable ( X ) is:
\[
P\{X > x\}
\]

\subsection*{Why do we care about tails?}

\begin{itemize}
\item Jobs delayed more than 24 hours
\item Packets dropped due to full buffer
\end{itemize}

\subsection*{Key Idea (from lecture)}

\begin{itemize}
\item Mean $\rightarrow$ average behavior (easy)
\item Tail $\rightarrow$ rare extreme events (hard)
\end{itemize}

\subsection*{Why are tails hard to compute?}

\begin{itemize}
\item Requires summing many small probabilities
\end{itemize}

\section*{Tail Example 1}

\subsection*{Problem}

Suppose distributing ( n ) jobs among ( n ) servers at random.
Find:
\[
P\{X \ge k\}
\]
where ( X ) = number of jobs assigned to a server.

\subsection*{Reasoning (Step-by-step from lecture)}

\begin{enumerate}
\item Jobs are assigned randomly
\item Most servers get few jobs
\item Occasionally, one server gets many jobs
\item This is a \textbf{tail event (rare overload)}
\end{enumerate}

\subsection*{Mathematical Formulation}

\[
X \sim \text{Binomial}(n, p)
\]

\[
P\{X \ge k\} = \sum_{i=k}^{n} \binom{n}{i} p^i (1-p)^{n-i}
\]

\section*{Tail Example 2 (Poisson Tail)}

\subsection*{Problem}

Jobs arrive according to a Poisson process with rate ( $\lambda$ ).
Find:
\[
P\{X \ge k\}
\]

\subsection*{Reasoning (Step-by-step from lecture)}

\begin{enumerate}
\item Jobs arrive randomly over time
\item Most hours have moderate arrivals
\item Sometimes there is a burst of arrivals
\item This is a \textbf{tail event (sudden spike in load)}
\end{enumerate}

\subsection*{Mathematical Formulation}

\[
X \sim \text{Poisson}(\lambda)
\]

\[
P\{X \ge k\} = \sum_{i=k}^{\infty} e^{-\lambda} \frac{\lambda^i}{i!}
\]

\section*{Why Tail Bounds?}

\subsection*{Observation}

Exact tail probabilities are often hard to compute.

\subsection*{Exact Computation}

\[
P\{X \ge k\} = \sum_{i=k}^{\infty} e^{-\lambda} \frac{\lambda^i}{i!}
\]

\[
P\{X \ge k\} = \sum_{i=k}^{n} \binom{n}{i} p^i (1-p)^{n-i}
\]

\begin{itemize}
\item These are \textbf{complicated, long sums}
\end{itemize}

\subsection*{Key Idea (from lecture)}

Instead of exact value:

\begin{itemize}
\item Find an \textbf{upper bound}
\item Easier to compute
\item Still gives a \textbf{guarantee}
\end{itemize}

\[
P(X \ge k) \le \text{simple bound}
\]

\textbf{Key idea:} Approximate safely instead of computing exactly

\section*{Running Example (Used Across Inequalities)}

\begin{itemize}
\item Experiment: Flip a fair coin ( n ) times
\item Random variable: ( $X =$ number of heads )
\item Distribution:
\[
X \sim \text{Binomial}(n, 1/2)
\]
\item Mean:
\[
E[X] = \frac{n}{2}
\]
\end{itemize}

\subsection*{Question}

\[
P\left(X \ge \frac{3}{4}n\right)
\]

\begin{itemize}
\item Tail region: unusually many heads
\end{itemize}

\section*{2. Markov’s Inequality}

\subsection*{Key Idea}

\begin{itemize}
\item Uses only the mean
\item Large values $\Rightarrow$ rare
\end{itemize}

\subsection*{Definition}

\[
P(X \ge a) \le \frac{E[X]}{a}
\]

\subsection*{Proof (Step-by-step as per lecture flow)}

\begin{enumerate}
\item Consider random variable ( $X \ge 0$ )
\item Split expectation:
\[
E[X] = \sum_x x \cdot P(X = x)
\]
\item Focus on values where ( $X \ge a$ ):
\[
E[X] \ge \sum_{x \ge a} x \cdot P(X = x)
\]
\item Since ( $x \ge a$ ):
\[
\ge \sum_{x \ge a} a \cdot P(X = x)
\]
\item Factor out ( $a$ ):
\[
= a \cdot P(X \ge a)
\]
\item Therefore:
\[
E[X] \ge a \cdot P(X \ge a)
\]
\item Rearranging:
\[
P(X \ge a) \le \frac{E[X]}{a}
\]
\end{enumerate}

\subsection*{Example (from lecture)}

\begin{itemize}
\item Avg = 10 $\Rightarrow$ few values can be 100+
\end{itemize}

\section*{3. Chebyshev’s Inequality}

\subsection*{Definition}

\[
P(|X - \mu| \ge k\sigma) \le \frac{1}{k^2}
\]

\subsection*{Proof (Step-by-step)}

\begin{enumerate}
\item Apply Markov inequality on:
\[
Y = (X - \mu)^2
\]
\item Then:
\[
P(|X - \mu| \ge k\sigma) = P((X - \mu)^2 \ge k^2\sigma^2)
\]
\item Apply Markov:
\[
\le \frac{E[(X - \mu)^2]}{k^2\sigma^2}
\]
\item Since:
\[
E[(X - \mu)^2] = \sigma^2
\]
\item Therefore:
\[
P(|X - \mu| \ge k\sigma) \le \frac{1}{k^2}
\]
\end{enumerate}

\subsection*{Example}

(As per lecture: deviation from mean bounded using variance)

\section*{4. Chernoff Bound and Poisson Tail}

\subsection*{Key Idea (Chernoff Bound)}

\begin{itemize}
\item Provides \textbf{tighter bounds} than Markov/Chebyshev
\item Works well for sums of independent random variables
\end{itemize}

\subsection*{Definition (Form Used in Lecture Context)}

For sum of independent variables:
\[
P(X \ge (1+\delta)\mu) \le \text{exponentially decreasing bound}
\]

\subsection*{Proof (Step-by-step structure from lecture)}

\begin{enumerate}
\item Start with:
\[
P(X \ge a)
\]
\item Transform using exponential:
\[
P(e^{tX} \ge e^{ta})
\]
\item Apply Markov inequality:
\[
\le \frac{E[e^{tX}]}{e^{ta}}
\]
\item Use independence:
\[
E[e^{tX}] = \prod E[e^{tX_i}]
\]
\item Optimize over ( $t$ )
\end{enumerate}

\subsection*{Poisson Tail}

\[
P(X \ge k) = \sum_{i=k}^{\infty} e^{-\lambda} \frac{\lambda^i}{i!}
\]

\begin{itemize}
\item Chernoff bound gives \textbf{simpler upper bound instead of sum}
\end{itemize}

\subsection*{Example}

Applied on:

\begin{itemize}
\item Binomial distribution (coin flips)
\item Poisson arrivals
\end{itemize}

\section*{5. Jensen’s Inequality}

\subsection*{Convex Functions}

A function ( $f$ ) is convex if:
\[
f(\lambda x + (1-\lambda)y) \le \lambda f(x) + (1-\lambda)f(y)
\]

\subsection*{Jensen’s Inequality}

\[
f(E[X]) \le E[f(X)]
\]

\subsection*{Proof (Step-by-step from lecture structure)}

\begin{enumerate}
\item Use convexity definition
\item Extend from two points to expectation
\item Apply linearity of expectation
\item Obtain:
\[
f(E[X]) \le E[f(X)]
\]
\end{enumerate}

\subsection*{Example}

Applying convex function on expectation vs expectation of function

\section*{6. Case Study: System Level Understanding}

\subsection*{Problem}

How to predict worst-case demand on server?

\subsection*{Reasoning (Step-by-step from lecture)}

\begin{enumerate}
\item Model jobs assigned randomly
\item Identify overload as tail event
\item Exact computation:
\begin{itemize}
\item Requires large summations
\end{itemize}
\item Use tail bounds:
\begin{itemize}
\item Markov $\rightarrow$ simple but loose
\item Chebyshev $\rightarrow$ uses variance
\item Chernoff $\rightarrow$ tight bound
\end{itemize}
\item Choose bound based on:
\begin{itemize}
\item Accuracy requirement
\item Computational simplicity
\end{itemize}
\end{enumerate}

\subsection*{Conclusion from Lecture}

\begin{itemize}
\item Tail bounds allow:
\begin{itemize}
\item Avoiding exact computation
\item Getting guaranteed upper bounds
\item Predicting rare extreme events
\end{itemize}
\end{itemize}

\section*{FINAL SUMMARY (FOR EXAM)}

\begin{itemize}
\item Tail:
\[
P(X > x)
\]
\item Exact tail:
$\rightarrow$ Hard (long sums)

\item Strategy:
$\rightarrow$ Use bounds
\end{itemize}

\begin{center}
\begin{tabular}{|c|c|c|}
\hline
Inequality & Uses & Strength \\
\hline
Markov & Mean & Weak \\
Chebyshev & Mean + Variance & Medium \\
Chernoff & Distribution structure & Strong \\
\hline
\end{tabular}
\end{center}

\begin{itemize}
\item Jensen:
\[
f(E[X]) \le E[f(X)]
\]
\item Application:
$\rightarrow$ Predict rare events (server overload, bursts)
\end{itemize}

\end{document}